\documentclass[../../main.tex]{subfiles}
\begin{document}

\begin{center}
  {\fontsize{48pt}{58pt}\selectfont \textbf{1985}\normalsize\textbf{年度}}
\end{center}

\begin{table}[htb]
  \centering
  \caption{出題テーマと難易度}
  \begin{tabular}{|c|p{6.2cm}|c|c|c|}
    \hline
    \textbf{問題番号} & \textbf{テーマ} & \textbf{思考量} & \textbf{計算量} & \textbf{総合} \\
    \hline
    第1問 & 内積と辺の長さを指定した三角形の存在条件 & 難 & 難 & 難 \\
    \hline
    第2問 & 一次変換で不変な三角形の頂点 &  &  &  \\
    \hline
    第3問 & 複素数の線形結合と二次方程式の零点 &  &  &  \\
    \hline
    第4問 & 球面と平面の交円および整数解 &  &  &  \\
    \hline
    第5問 & 2枚の硬貨の表の枚数の分布 &  &  &  \\
    \hline
    第6問 & 積分方程式を満たす連続関数 &  &  &  \\
    \hline
  \end{tabular}
\end{table}

\begin{table}[htb]
  \centering
  \caption{解答}
  \renewcommand{\arraystretch}{1.4}
  \begin{tabular}{|c|c|p{9cm}|}
    \hline
    \textbf{問題番号} & \textbf{小問} & \textbf{答え} \\
    \hline
    第1問 & -- & $p> -\dfrac{q^2}{4}$ \\
    \hline
    第2問 & -- &  \\
    \hline
    第3問 & (i) & 任意の複素数を$\alpha,\beta$の実係数の一次結合で表せる \\
    \cline{2-3}
    & (ii) & $D$の条件と図は解答を参照 \\
    \hline
    第4問 & (1) & $PQ\le2\sqrt{\dfrac{2}{3}}$ \\
    \cline{2-3}
    & (2) & $r=3k+1$なら$(k,k,k+1)$の順列，$r=3k+2$なら$(k,k+1,k+1)$の順列．$3$の倍数なら解なし \\
    \hline
    第5問 & (1) & $(a,b)=(p,p)$の1組 \\
    \cline{2-3}
    & (2) & 存在しない \\
    \hline
    第6問 & -- & $a=0,b\ge1$なら$f(x)=Cx^{b-1}$（$C$は定数）．それ以外は$f(x)=0$ \\
    \hline
  \end{tabular}
\end{table}

\paragraph{試験形態}
6問．試験時間は要確認．

\end{document}
